\begin{lemma}
\label{lemma-cone-koszul-abstract}
Let $R$ be a ring. Let $\varphi : E \to R$ be an $R$-module map.
Let $f \in R$. Set $E' = E \oplus R$ and define $\varphi' : E' \to R$
by $\varphi$ on $E$ and multiplication by $f$ on $R$.
The complex $K_\bullet(\varphi')$ is isomorphic to the
cone of the map of complexes
$$
f :
K_\bullet(\varphi)
\longrightarrow
K_\bullet(\varphi).
$$
\end{lemma}

\begin{proof}
Denote $e_0 \in E'$ the element $1 \in R \subset R \oplus E$.
By our definition of the cone above we see that
$$
C(f)_n = K_n(\varphi) \oplus K_{n - 1}(\varphi) =
\wedge^n(E) \oplus \wedge^{n - 1}(E) = \wedge^n(E')
$$
where in the last $=$ we map $(0, e_1 \wedge \ldots \wedge e_{n - 1})$
to $e_0 \wedge e_1 \wedge \ldots \wedge e_{n - 1}$ in $\wedge^n(E')$.
A computation shows that this isomorphism is compatible with
differentials. Namely, this is clear for elements of the first
summand as $\varphi'|_E = \varphi$ and $d_{C(f)}$ restricted to
the first summand is just $d_{K_\bullet(\varphi)}$.
On the other hand, if $e_1 \wedge \ldots \wedge e_{n - 1}$
is in the second summand, then
$$
d_{C(f)}(0, e_1 \wedge \ldots \wedge e_{n - 1}) =
fe_1 \wedge \ldots \wedge e_{n - 1}
- d_{K_\bullet(\varphi)}(e_1 \wedge \ldots \wedge e_{n - 1})
$$
and on the other hand
\begin{align*}
& d_{K_\bullet(\varphi')}(0, e_0 \wedge e_1 \wedge \ldots \wedge e_{n - 1}) \\
& =
\sum\nolimits_{i = 0, \ldots, n - 1}
(-1)^i \varphi'(e_i)e_0 \wedge \ldots \wedge \widehat{e_i}
\wedge \ldots \wedge e_{n - 1} \\
& =
fe_1 \wedge \ldots \wedge e_{n - 1} +
\sum\nolimits_{i = 1, \ldots, n - 1}
(-1)^i \varphi(e_i)e_0 \wedge \ldots \wedge \widehat{e_i}
\wedge \ldots \wedge e_{n - 1} \\
& =
fe_1 \wedge \ldots \wedge e_{n - 1} -
e_0 \left(\sum\nolimits_{i = 1, \ldots, n - 1}
(-1)^{i + 1} \varphi(e_i)e_1 \wedge \ldots \wedge \widehat{e_i}
\wedge \ldots \wedge e_{n - 1}\right)
\end{align*}
which is the image of the result of the previous computation.
\end{proof}
